\chapter[Comparative Geometry]{Comparative Geometry: UOD, Fisher, and Classical Divergences}\label{ch:sec:comparative}


This chapter situates the UOD relative to existing geometric frameworks---the Fisher metric, the Hellinger distance, Pearson's $\chi^2$ divergence. These comparisons clarify what is new about the WIS geometry and prepare the ground for the discussion of related work in Chapter~\ref{ch:sec:discussion}.


\section{Introduction}\label{ch:sec:introduction-comparative_geometry}

The UOD is not an isolated geometric quantity.  This chapter compares
it with the Fisher information metric (the classical Riemannian metric
of information geometry) and with standard statistical divergences,
both locally and globally.

\section{Fisher Information Geometry}\label{ch:sec:fisher-information-geometry}

Throughout this chapter let $(\mathcal P,g)$ denote the posterior manifold equipped with the pullback Riemannian metric introduced in Chapter~\ref{ch:sec:pullback-riemannian-metric}.
The purpose of this chapter is to compare the geometric structure
induced by Weighted Incidence Structures with the classical Fisher
information geometry.
\begin{remark}
The results of this chapter establish structural
relationships. They do not assert that the pullback metric always
coincides with the Fisher metric. Additional assumptions may be
required for such an identification.
\end{remark}
\paragraph{Fisher Information Metric.}
Let $\Theta$ be a smooth statistical model with probability distributions
$p(x;\theta).$
The Fisher information metric is defined by
\[
g_F(\theta)
=
\mathbb E_\theta
\!\left[
\nabla_\theta\log p(X;\theta)
\nabla_\theta\log p(X;\theta)^T
\right].
\]
It is a positive-semidefinite symmetric tensor and, under regularity
assumptions, defines a Riemannian metric on the statistical manifold.

\section{Comparison of the Two Geometries}\label{ch:sec:comparison-of-the-two-geometries}
The pullback metric constructed in Chapter~\ref{ch:sec:pullback-riemannian-metric} is
$g = F^{*}g_E,$
where
$F = Q\circ L$
is the canonical logarithmic coordinate map.
Unlike the Fisher metric, which is obtained from expectations of score
functions, the pullback metric is induced by the differential of a
globally defined logarithmic embedding modulo gauge equivalence.
Thus the two constructions originate from different geometric
principles.
\begin{proposition}[Shared Properties]
\label{ch:prop:10-1}
Both the Fisher metric and the pullback metric are symmetric
positive-definite tensor fields whenever they are defined on regular
statistical models.
\end{proposition}
\begin{proof}
The Fisher metric is symmetric by construction and positive definite
under the usual regularity assumptions.
The pullback metric is symmetric and positive definite by Theorem~\ref{ch:thm:7-6}.
\end{proof}

\begin{proposition}[Coordinate Invariance]
\label{ch:prop:10-2}
Both metrics are invariant under smooth reparameterizations of the
underlying manifold.
\end{proposition}
\begin{proof}
The Fisher metric is invariant under smooth coordinate transformations
by classical information geometry.
The pullback metric is defined as the pullback of the Euclidean metric
by a diffeomorphism and is therefore intrinsically
coordinate-independent. 
\end{proof}
\subsection{Structural Differences}\label{ch:sec:structural-differences}
The Fisher metric depends on
\begin{itemize}
\item the statistical model,
\item the likelihood function,
\item expectations with respect to the model.
\end{itemize}
The pullback metric developed in this work depends on
\begin{itemize}
\item posterior distributions,
\item logarithmic coordinates,
\item quotient geometry induced by additive gauge symmetry.
\end{itemize}
Consequently, the pullback metric remains meaningful even in situations
where no differentiable statistical parameterization is explicitly
available.
\begin{theorem}[Intrinsic Definition]\label{ch:thm:10-3}

The pullback metric is intrinsically defined on the posterior manifold
independently of any statistical parameterization.
\end{theorem}
\begin{proof}
By Theorem~\ref{ch:thm:5-4}, the canonical logarithmic coordinate map
is a local diffeomorphism.
The pullback metric is therefore defined entirely through
the logarithmic embedding, which depends only on the posterior manifold and the Euclidean structure
of the target space.
No external parameterization is required. 
\end{proof}

\begin{corollary}[Generality]
\label{ch:cor:10-4}
The differential geometry induced by Weighted Incidence Structures
extends beyond the classical setting of parameterized statistical
models.
\end{corollary}
\begin{proof}
Immediate from Theorem~\ref{ch:thm:10-3}. 
\end{proof}

\section{Discussion}\label{ch:sec:discussion-comparative_geometry}
The Fisher metric and the pullback metric share several fundamental
geometric properties:
\begin{itemize}
\item symmetry;
\item positive definiteness;
\item coordinate invariance;
\item smooth dependence on the underlying manifold.
\end{itemize}
However, they arise from distinct constructions.
The Fisher metric is expectation-based and model-dependent.
The pullback metric introduced in this appendix is induced by the
logarithmic geometry of posterior distributions modulo gauge
transformations.
This distinction suggests that the present framework should be viewed as
complementary to classical Information Geometry rather than as a
replacement for it.
Future work may identify conditions under which the two metrics coincide
or become naturally isometric.

\begin{theorem}[Local Quadratic Expansion]\label{ch:thm:25-1}

For every \(P\in\mathcal P^\circ\) and every
\(\delta\in T_P\mathcal P^\circ\),
\[ \mathcal D(P+\delta,P) =
g_P(\delta,\delta) + o(\|\delta\|^2). \]
\end{theorem}
\begin{proof}
Taylor expansion of \(F\) gives
$F(P+\delta)=F(P)+DF_P(\delta)+o(\|\delta\|),$
therefore
\[ \mathcal D(P+\delta,P) =
\|DF_P(\delta)\|^2 + o(\|\delta\|^2). \]
Since \(g=F^{*}g_E\), the conclusion follows. 
\end{proof}

\paragraph{Running example: differential privacy.}
The comparison between the pullback metric and the Fisher metric has a concrete interpretation for differential privacy (Section~\ref{ch:sec:differential-privacy-and-database-queries}).  The Laplace mechanism achieves $\mathcal D=6/\varepsilon^2$, and Theorem~\ref{ch:thm:25-1} shows that locally the UOD is quadratic in the Fisher metric.  Hence for small privacy budgets $\varepsilon$ the loss of utility is accurately captured by the Fisher geometry---a fact that Bayesian privacy analysts can exploit directly.

\section{Explicit Formula}\label{ch:sec:explicit-formula}
\begin{proposition}[Pullback Metric Formula]
\label{ch:prop:25-2}
For tangent vectors \(u,v\),
\[ g(u,v) = \sum_i\frac{u_iv_i}{P_i^2} -
\frac1n
\left(\sum_i\frac{u_i}{P_i}\right)
\left(\sum_i\frac{v_i}{P_i}\right). \]
\end{proposition}
\begin{proof}
Since
\[ DF_P=Q\circ DL_P, \qquad
DL_P(u)=\left(\frac{u_i}{P_i}\right), \]
and \(Q\) is orthogonal,
\[ g(u,v) = \langle Q\,DL_P u,Q\,DL_P v\rangle =
\langle DL_P u,Q\,DL_P v\rangle. \]
Expanding \(Q=I-\frac1n\mathbf{11}^T\) yields the formula. 
\end{proof}
\paragraph{Fisher Metric.}
The Fisher information metric is
$g_F(u,v) = \sum_i\frac{u_iv_i}{P_i}.$

\paragraph{Matrix Representation.}
Let
\[
D(P)=\operatorname{diag}\!\left(\frac1{P_i^2}\right),
\qquad v(P)=\left(\frac1{P_i}\right).
\]
Then
$G(P) = D(P)-\frac1n v(P)v(P)^T$
is the matrix representing the pullback metric.
The Fisher metric is represented by
\[
G_F(P)=\operatorname{diag}\!\left(\frac1{P_i}\right).
\]

\section{Uniform Equivalence}\label{ch:sec:uniform-equivalence}
\begin{theorem}[Uniform Metric Equivalence]
\label{ch:thm:25-3}
Let \(K\subset\mathcal P^\circ\) be compact.
Then there exist constants
$0<c_K\le C_K<\infty$
such that
$c_K\,g_F(u,u) \le g(u,u) \le C_K\,g_F(u,u)$
for every \(P\in K\) and every tangent vector
\(u\in T_P\mathcal P^\circ\).
\end{theorem}
\begin{proof}
Consider the normalized symmetric matrix
$A(P) = G_F(P)^{-1/2} G(P) G_F(P)^{-1/2}.$
Since both metrics are symmetric,
$g(u,u) = x^T A(P) x,$
where
$x = G_F(P)^{1/2}u.$
The matrix \(A(P)\) is symmetric and positive definite on the tangent
space. Hence, by the Rayleigh quotient theorem,
\[ \lambda_{\min}(A(P)) \|x\|^2 \le x^T A(P) x
\le \lambda_{\max}(A(P)) \|x\|^2. \]
Because
$\|x\|^2 = g_F(u,u),$
we obtain
\[ \lambda_{\min}(A(P))\,g_F(u,u) \le g(u,u)
\le \lambda_{\max}(A(P))\,g_F(u,u). \]
The entries of \(A(P)\) depend continuously on \(P\). Since \(K\) is compact,
its minimum and maximum eigenvalues are attained:
\[ c_K =
\min_{P\in K}\lambda_{\min}(A(P)),
\qquad C_K =
\max_{P\in K}\lambda_{\max}(A(P)). \]
Because \(A(P)\) is positive definite on the tangent space,
$0<c_K\le C_K<\infty.$
Therefore
$c_K\,g_F \le g \le C_K\,g_F.$
Consequently, the two metrics induce the same local topology on every
compact subset of the interior simplex. 
\end{proof}
\paragraph{Consequences.}
\begin{corollary}[Local Topological Equivalence]
\label{ch:cor:25-4}
On every compact subset:
\begin{itemize}
\item the pullback and Fisher metrics generate the same local topology;
\item they define equivalent notions of convergence;
\item infinitesimal estimates can be transferred between the two
    geometries up to uniform constants.
\end{itemize}
\end{corollary}

\paragraph{Interpretation.}
The Fisher metric arises from the Hessian of relative entropy, whereas
the WIS metric is obtained by transporting Euclidean geometry through
logarithmic quotient coordinates.
Although their global constructions differ, Theorem~\ref{ch:thm:25-3} shows that they
are locally equivalent on compact subsets of the posterior manifold.

\paragraph{Outlook.}
Having established the local relationship with Fisher geometry, we now
turn to global comparisons between the Universal Optimality Defect and
classical statistical divergences.
\section{Comparison with Total Variation}\label{ch:sec:comparison-with-total-variation}
The Total Variation distance is
$TV(P,Q)=\frac12\sum_{i}|P_i-Q_i|.$
\begin{theorem}
\label{ch:thm:26-1}
There exist sequences of probability distributions for which the Total
Variation distance remains bounded while the UOD diverges to infinity.
\textbf{Proof.}
Let
$Q=\left(\frac13,\frac13,\frac13\right),$
and
\[ P_\varepsilon= \left( \varepsilon,
\frac{1-\varepsilon}{2}, \frac{1-\varepsilon}{2}
\right). \]
Then
\[ TV(P_\varepsilon,Q)\longrightarrow\frac13,
\]
whereas
$\log\frac{\varepsilon}{1/3}\longrightarrow-\infty.$
After projection onto the quotient hyperplane the logarithmic coordinates still diverge, and by Theorem~\ref{ch:thm:26-1} the UOD
$\mathcal{D}(P_\varepsilon,Q)\longrightarrow\infty.$
Hence no bounded function of Total Variation can globally characterize
the UOD. 
\end{theorem}

\paragraph{Interpretation.}
The UOD is not equivalent to any classical divergence: unlike the Fisher metric, it is flat; unlike the Hellinger distance, it is unbounded; unlike $\chi^2$, it is quadratic in log-ratios rather than in ratios.  These differences are not defects but signatures of the fact that the UOD measures distance from a structural condition, not between two distributions.

\subsection{Comparison with Hellinger Distance}\label{ch:thm:26-2}
The squared Hellinger distance is
\[ H^{2}(P,Q)= \frac12
\sum_{i}(\sqrt{P_i}-\sqrt{Q_i})^{2}. \]
\begin{theorem}

The Hellinger distance is bounded whereas the UOD is unbounded.
\textbf{Proof.}
Using
$H^{2}(P,Q) = 1-\sum_{i}\sqrt{P_iQ_i},$
one immediately obtains $0\le H(P,Q)\le1.$
By Theorem~\ref{ch:thm:26-1}, the UOD diverges along
the boundary of the simplex. Therefore the UOD cannot be represented
globally as a function of the Hellinger distance. 
\end{theorem}
\subsection{Local Comparison with Pearson's Chi-Squared Divergence}\label{ch:sec:local-comparison-with-pearsons-chi-squared-divergence}
Pearson's divergence is
\[ \chi^{2}(P,Q) = \sum_{i}
\frac{(P_i-Q_i)^2}{Q_i}. \]
Let
\[ P=Q+\delta, \qquad
\sum_{i}\delta_{i}=0. \]
Using
$\log(1+x)=x-\frac{x^2}{2}+O(x^{3}),$
one obtains
\[ \chi^{2}(P,Q) = \sum_{i}
\frac{\delta_i^2}{Q_i} + O(\|\delta\|^{3}), \]
whereas
\[ \mathcal{D}(P,Q) = \sum_{i}
\frac{\delta_i^2}{Q_i^2} - \frac1n \left(
\sum_{i}\frac{\delta_i}{Q_i} \right)^{2} +
O(\|\delta\|^{3}). \]
Thus both quantities possess quadratic local expansions, but they induce
different quadratic forms.
\section{First Structural Comparison with KL}\label{ch:sec:first-structural-comparison-with-kl}
The Kullback--Leibler divergence satisfies
\[ D_{KL}(P\|Q) = \sum_{i} P_i
\log\frac{P_i}{Q_i}. \]
Introducing the logarithmic field
$a_i=\log\frac{P_i}{Q_i},$
we obtain the identity
\[ \boxed{
D_{KL}(P\|Q)=E_P[a].
} \]
By Corollary~\ref{ch:cor:21-3},
\[ \boxed{
\mathcal{D}(P,Q)
=
n\,Var_U(a),
} \]
where (U) denotes the uniform distribution.
Hence the KL divergence and the UOD are both functionals of the same
log-likelihood field, but evaluated in fundamentally different ways: the
KL is its expectation with respect to (P), whereas the UOD is its
centered second moment with respect to the uniform measure.
This observation motivates the functional framework developed in the
next chapter.
\paragraph{Summary.}
The rigorous comparison established in this chapter shows:
\begin{itemize}
  \item the UOD is locally equivalent to Fisher geometry;
  \item it is globally distinct from Total Variation;
  \item it is globally distinct from Hellinger distance;
\end{itemize}
